\documentclass[11pt,a4paper]{article}

\usepackage[margin=1in]{geometry}
\usepackage{amsmath,amssymb}
\usepackage{hyperref}
\usepackage{booktabs}
\usepackage{enumitem}
\usepackage{microtype}
\usepackage{parskip}

\hypersetup{
  colorlinks=true,
  linkcolor=black,
  citecolor=black,
  urlcolor=blue
}

\title{\textbf{LLM-Tor: Unlinkable Access to Commercial LLM APIs via Blind Signatures and Tor}}
\author{Prince Gupta}
\date{March 2026}

\begin{document}

\maketitle

\begin{abstract}
LLM-Tor is a cryptographically enforced privacy layer that enables unlinkable access to commercial
Large Language Model (LLM) APIs. It separates payment identity from model usage using RSA blind
signatures (RFC 9474) and onion-routed communication via Tor. The system prevents correlation between
user identity and chat content at the proxy layer while maintaining compatibility with frontier commercial
LLM providers.
\end{abstract}

\tableofcontents
\newpage

%% ─────────────────────────────────────────────────────────────────────────────
\section{Introduction}

Commercial LLM APIs require authenticated access tied to user accounts. This creates a structural privacy
issue: prompts and responses become permanently associated with identifiable users.

LLM-Tor introduces a cryptographic separation between:
\begin{itemize}
  \item Payment identity
  \item Usage authorization
  \item Model inference requests
\end{itemize}

The goal is to ensure that even the proxy operator cannot link a specific prompt to a specific paying
customer.

%% ─────────────────────────────────────────────────────────────────────────────
\section{Background}

\subsection{Blind Signatures}

LLM-Tor uses RSA Blind Signatures as defined in RFC 9474, specifically the
\texttt{RSABSSA\allowbreak-SHA384\allowbreak-PSS\allowbreak-Randomized} suite.

A blind signature scheme allows a signer to sign a message without learning its contents. The protocol
consists of:
\begin{enumerate}
  \item Blinding
  \item Signing
  \item Unblinding
  \item Verification
\end{enumerate}

Security relies on the hardness of RSA and the one-more RSA assumption.

Let $m$ be a random token. The client computes a blinding factor $r$ and sends:
\[
  m' = m \cdot r^e \pmod{N}
\]
The server blind-signs:
\[
  s' = (m')^d \pmod{N}
\]
The client unblinds:
\[
  s = s' \cdot r^{-1} \pmod{N}
\]
The resulting signature $s$ is valid for $m$, but the signer never sees $m$.

\subsection{Tor Onion Routing}

Tor provides layered encryption and multi-hop routing to hide client IP addresses. LLM-Tor routes all
anonymous redemption traffic through Tor. In future, the inference endpoint will be exclusively accessible
as a Tor onion service.

Reasons for using Tor:
\begin{itemize}
  \item Prevent IP-level linkage during token redemption
  \item Eliminate DNS metadata leakage
  \item Provide probabilistic anonymity against network observers
\end{itemize}

LLM-Tor does not rely on Tor for cryptographic unlinkability; Tor and blind signatures are independent,
complementary protections.

%% ─────────────────────────────────────────────────────────────────────────────
\section{System Architecture}

\subsection{Components}

\begin{itemize}
  \item Payment Service
  \item Credit Ledger
  \item Blind Signing Service
  \item Token Redemption Proxy
  \item Moderation Layer
  \item Upstream LLM Provider
\end{itemize}

\subsection{Protocol Phases}

\subsubsection{Phase 1: Credit Purchase}

\begin{enumerate}
  \item User purchases credits via standard payment provider.
  \item Server records credit balance linked to user identity.
\end{enumerate}
Identity exists in this phase.

\subsubsection{Phase 2: Blind Token Issuance}

\begin{enumerate}
  \item Client generates random token $T$.
  \item Client blinds $T$ using model-specific public key.
  \item Server verifies credit availability.
  \item Server blind-signs the blinded token.
  \item Client unblinds signature.
\end{enumerate}
The server never observes $T$ in unblinded form. Identity exists in this phase.

\subsubsection{Phase 3: Anonymous Redemption}

\begin{enumerate}
  \item Client establishes Tor circuit.
  \item Client submits $(T, \sigma)$.
  \item Server verifies RSA signature.
  \item Server checks token not previously spent.
  \item Server forwards prompt to upstream LLM.
\end{enumerate}
No identity is present in this phase.

%% ─────────────────────────────────────────────────────────────────────────────
\section{Token Structure}

Each token contains:
\begin{itemize}
  \item A random 128-bit value.
  \item Model identifier: implicitly linked by choosing the public key specific to that model.
  \item Key identifier (for rotation) [planned].
\end{itemize}
Spent tokens are stored as $\mathrm{hash}(T)$ to prevent replay.

%% ─────────────────────────────────────────────────────────────────────────────
\section{Abuse Prevention System}
\label{sec:abuse}

\subsection{Motivation}

Per-request LLM tokens are single-use and fully unlinkable; there is no stable identifier that the proxy
can use to rate-limit or ban a misbehaving user. Content moderation rejects individual bad requests, but
does not prevent a determined adversary from repeatedly issuing new tokens and continuing to submit
harmful content. The abuse-token scheme provides the proxy with a \emph{stable but anonymous} handle
per user. This is largely motivated from legal compliance and ability to adhere to upstream TOS.

\subsection{Dual-Token Design}

Two tokens serve complementary roles:

\begin{center}
\renewcommand{\arraystretch}{1.4}
\begin{tabular}{@{}lll@{}}
\toprule
\textbf{Token} & \textbf{Structure} & \textbf{Validity} \\
\midrule
Permanent ($A_p$) & 48 random bytes & Lifetime of account \\
Transient ($A_t$) & 44 random bytes $\|$ 4-byte epoch & One calendar month \\
\bottomrule
\end{tabular}
\end{center}

\begin{description}
  \item[Permanent token $A_p$.]  Issued once per account. Remains valid indefinitely. Blacklisting
    $A_p$ permanently bans the anonymous pseudonym.
  \item[Transient token $A_t$.]  Issued once per calendar month. Blacklisting $A_t$ provides a
    time-bounded ban; the permanent token continues to provide a long-lived~identifier.
\end{description}

The epoch is defined as $\text{epoch} = \text{year} \times 12 + \text{month}$ (month 1-indexed), encoded
as a big-endian 32-bit unsigned integer appended to the random bytes. Each token type uses a dedicated
RSA key pair, independent of the per-model keys.

Both tokens are blind-signed, so the server cannot correlate issuance (where the user is authenticated)
with redemption (where the user is anonymous over Tor).

\subsection{Issuance}
\label{sec:abuse:issuance}

Issuance follows the same blind-signature protocol as per-request tokens, using two dedicated RSA key
pairs:

\begin{enumerate}
  \item Client generates random bytes (48 for permanent; 44 random $\|$ 4-byte epoch for transient).
  \item Client blinds with the corresponding abuse public key and sends the blinded token to the server
        over an authenticated session.
  \item Server checks that the token type has not already been issued for this account in the current
        epoch:
        \begin{itemize}
          \item Permanent: rejects if \texttt{PermanentAbuseTokenIssuedAt} is non-nil.
          \item Transient: rejects if \texttt{TransientAbuseTokenEpoch} equals the current epoch.
        \end{itemize}
  \item Server blind-signs and records the issuance flag atomically (protected by a per-user distributed
        semaphore to prevent concurrent double-issuance).
  \item Client unblinds and verifies the signature.
\end{enumerate}

\subsection{Blacklisting Mechanism}

Blacklisted tokens are stored by their hash:
\[
  \text{DocID} = \mathrm{base64}\!\left(\mathrm{SHA\text{-}256}(A)\right)
\]

This mirrors the pattern used for spent per-request tokens. The server checks both the permanent and
transient blacklists on every anonymous request before forwarding to the upstream LLM.

\subsection{Anonymity Properties}

\paragraph{Unlinkability.}
Because both $A_p$ and $A_t$ are blind-signed, the server cannot link the issuance event (authenticated
session) to any specific anonymous request. The permanent token creates a persistent \emph{anonymous
pseudonym}: all requests bearing the same $A_p$ are linkable to each other, but not to the paying
identity.

\paragraph{Pseudonymity trade-off.}
The scheme intentionally trades full per-request unlinkability for pseudonymity within a session-less
context. This is a deliberate design decision: without any stable identifier the proxy has no mechanism to
enforce abuse bans. The permanent token carries no real-world identity; it is simply a stable random
value whose signature the proxy can verify.

\paragraph{Tor independence.}
Abuse-token verification occurs at the application layer. Tor provides network-layer anonymity
independently; the two protections are orthogonal.

\subsection{Client-Side Backup and Recovery}

The client encrypts the stored tokens for backup using:
\begin{itemize}
  \item Key derivation: PBKDF2 (SHA-256, 600\,000 iterations, 128-bit random salt).
  \item Encryption: AES-256-GCM with a 96-bit random IV.
  \item Format: $\mathrm{base64}(\text{salt} \,\|\, \text{IV} \,\|\, \text{ciphertext})$.
\end{itemize}

The encrypted backup is always exported as a local file (\texttt{.llmtorbak}); server-side backup
storage is opt-in. In the server-sync case the server stores only the client-encrypted ciphertext and
cannot decrypt it.

%% ─────────────────────────────────────────────────────────────────────────────
\section{Security Properties}

\subsection{Unlinkability}

Under the security assumptions of RSA blind signatures, issued tokens cannot be linked to redeemed tokens.

\subsection{Single-Use Enforcement}

Tokens are marked spent atomically upon first successful verification. Replay attempts are rejected.

\subsection{Identity Separation}

Payment identity never appears in redemption requests. Redemption occurs exclusively over Tor.

%% ─────────────────────────────────────────────────────────────────────────────
\section{Moderation Requirement}

LLM-Tor performs content moderation prior to forwarding requests upstream.

Reasons:
\begin{itemize}
  \item Legal compliance in certain jurisdictions
  \item Abuse mitigation (spam, illegal content)
  \item Prevention of infrastructure misuse
\end{itemize}

Moderation is performed at the point where content arrives at the proxy; at that point the content is not
attributable to any user identity.

%% ─────────────────────────────────────────────────────────────────────────────
\section{Threat Model}

\subsection{Adversaries}

\begin{itemize}
  \item Honest-but-curious proxy operator
  \item External network observer
  \item Malicious user attempting double-spend
  \item Upstream LLM provider
\end{itemize}

\subsection{Security Goals}

\begin{itemize}
  \item Prevent proxy from linking prompt to payment identity
  \item Prevent token double spending
  \item Prevent IP-level identity leakage
\end{itemize}

\subsection{Out of Scope}

\begin{itemize}
  \item Global passive adversary controlling Tor
  \item Endpoint compromise
  \item Logging by upstream LLM providers
  \item Stylometric fingerprinting
\end{itemize}

%% ─────────────────────────────────────────────────────────────────────────────
\section{Key Rotation Policy}

TODO

%% ─────────────────────────────────────────────────────────────────────────────
\section{Limitations}

LLM-Tor does not provide absolute anonymity. It provides unlinkability at the proxy layer under standard
cryptographic assumptions, with the added benefit of network anonymity via Tor built into the official
client.

%% ─────────────────────────────────────────────────────────────────────────────
\section{Future Work}

\begin{itemize}
  \item Exclusive Tor onion hosting for the proxy microservice
  \item Key rotation
  \item Zero-knowledge token systems
  \item Anonymous payment integration
  \item Post-quantum blind signatures
  \item Decentralized moderation
\end{itemize}

%% ─────────────────────────────────────────────────────────────────────────────
\section{Conclusion}

LLM-Tor demonstrates that cryptographic blind signatures combined with onion routing can provide practical
unlinkability for commercial LLM API access, without requiring local model hosting.

\end{document}
